\chapter*{Preface}
\addcontentsline{toc}{chapter}{Preface}
\markboth{Preface}{Preface}

This book is an introduction to SpherePOP, a notation for writing computations as nested circles. Every circle names an operator and encloses its operands, and a computation proceeds by popping circles whose contents are ready, each popped circle being replaced by the value that its operator produces. The notation is small enough to be specified completely in a few pages, and yet it is rich enough to carry arithmetic, symbolic expressions, operations on pictures, and systems of rewriting rules. For that reason it makes a good setting in which to learn several ideas that recur throughout the study of programming languages: the difference between syntax and evaluation, the role of explicit grouping, the properties that make an evaluation order irrelevant, the way a system of sorts separates meaningful terms from meaningless ones, and the point at which rewriting becomes difficult.

The book is written for readers who know ordinary arithmetic notation, who have met proof by induction, and who are willing to follow a definition through a few examples before trusting it. No prior knowledge of lambda calculus, term rewriting, or type theory is assumed, although readers who have met these subjects will recognize familiar structures at several points and are told where to look for the standard treatments. Programming experience is helpful in the last chapter, where a complete reference interpreter of under three hundred lines is presented and its output is reproduced.

\paragraph{Organization.}
The first chapter motivates explicit grouping and introduces the circle drawings informally. Chapters 2 and 3 give the formal syntax and the single reduction rule. Chapter 4 proves the central structural results, namely that every computation terminates and that the order in which circles are popped does not affect the outcome. Chapter 5 introduces sorts and shows exactly where a well-sorted term can still fail to compute. Chapter 6 treats symbols and substitution, which allow a term to be partly evaluated before all of its inputs are known. Chapter 7 shows how the same machinery handles pictures and colours. Chapter 8 lets operators return circles instead of values, turning the notation into a rewriting system, and uses the classical criteria of termination and critical pairs to decide when the result is well behaved. Chapter 9 compares the notation with several neighbouring formalisms and indicates directions for extension. Chapter 10 presents the interpreter. The appendices contain solutions to selected exercises and a summary of notation.

\paragraph{Conventions.}
A circle is drawn as a closed curve, but in running text it is written $\cir{o\; t_1 \cdots t_n}$, with the operator name first and the operands following. The two forms are interchangeable. Operator names are set in sans-serif type, for example $\op{add}$ and $\op{rotate}$. Definitions, theorems, and examples are numbered within each chapter, and exercises are numbered by chapter. Exercises whose number is followed by a dagger have a solution in Appendix~A. Each proof ends with a square, and each example ends with a small triangle.

\paragraph{Using the interpreter.}
All numerical examples in the text have been computed with the interpreter of Chapter~10, and the transcript of a complete test run is reproduced there. Readers are encouraged to run the program, to change the examples, and to attempt the extensions suggested in the exercises at the end of that chapter.
